\section{Threats to Validity} \label{sec:threats}

\subsection{Internal Validity}

\textbf{Training-monitoring boundary effects.}
The transition from baseline training to the monitoring phase introduces a potential confound: the first few messages after training may receive artificially high anomaly scores simply because the sliding-window statistics are still stabilising.
We mitigate this by reporting metrics only from the monitoring phase and by using a sufficiently large moving-average window (100 messages) to dampen transient effects.

\textbf{Threshold selection.}
Defence effectiveness depends on chosen thresholds ($\tau_M = 0.05$, $\tau_A = 0.5$).
Different threshold values would produce different block rates and false-positive rates.
We report results for a single, documented configuration. Future work should conduct sensitivity analysis across threshold ranges.

\textbf{Measurement timing.}
Attack durations and inter-arrival times depend on Python thread scheduling and host system load.
Measured durations may deviate slightly from configured values (e.g., 5.38\,s observed vs.\ 5.0\,s configured).
We account for this by using actual wall-clock timestamps from logs rather than configured values in all metric computations.

\subsection{External Validity}

\textbf{Simulated vs.\ real traffic.}
The MAVLink SITL simulator produces a simplified traffic profile (four message types at regular intervals).
Real autonomous platforms exhibit more complex patterns including bursts during manoeuvres, variable-rate GPS updates, and multi-component communication.
Results may not directly transfer to physical deployments without validation on real hardware.

\textbf{Attack realism.}
The implemented attacks send well-formed MAVLink packets at elevated rates from known rogue source IDs.
Sophisticated attackers may use more subtle injection strategies (e.g., gradual drift, authenticated channels, adversarial sensor spoofing~\cite{sun2020lidar}, or adversarial examples~\cite{goodfellow2015adversarial, he2023advml} targeting ML models).
Our attack catalogue represents a starting point, and additional threat models should be evaluated.

\textbf{Single protocol.}
All experiments use MAVLink over UDP.
Autonomous systems using other protocols (DDS, ROS2, CAN bus) may exhibit different vulnerability profiles and require protocol-specific feature engineering.

\subsection{Construct Validity}

\textbf{Metric definitions.}
Our ``mission impact ratio'' (fraction of monitoring time with accuracy $<$ 0.9) is a proxy for true mission degradation.
A more precise metric would measure physical outcomes (trajectory deviation, task completion time), which requires integration with a physics-based simulator.

\textbf{Moving-average window size.}
The choice of window size (100 messages for static, 160 for randomised) affects reported accuracy values.
Smaller windows respond faster to attacks but exhibit more noise, while larger windows smooth out short bursts.
We report the window size explicitly to enable fair comparison with future work using different parameters.
